\section{Signal definitions and pre-registered signs}\label{app:signals}

The eight tanker signals used in Section~\ref{sec:partb}, the aggregation applied
when collapsing per-terminal or per-zone bands to one daily series, and the sign
committed before any fit. The sign convention: the spread $\text{HH}-\text{TTF}$ is
negative in nearly every period, so ``the spread narrows'' means it rises toward
zero --- a positive change.

\begin{center}\small
\begin{tabular}{p{3.4cm}p{1.6cm}p{7.4cm}c}
\toprule
signal & aggregation & mechanism & prior \\
\midrule
gas in transit, EU-bound & sum over EU zones & more cargo landing in Europe soon $\Rightarrow$ TTF falls & $+$ \\
gas in transit, unresolved & sum & destination unknown; no directional prior & none \\
US loading rate & sum over US berths & gas leaves the US balance (HH firms) and reaches Europe (TTF falls) & $+$ \\
EU discharge rate & sum over EU berths & Europe absorbing gas now $\Rightarrow$ TTF falls & $+$ \\
ballast returning to US & sum & tonnage returning to load $\Rightarrow$ future exports up & $+$ \\
laden voyage age & $n$-weighted mean & cargo floating rather than landing $\Rightarrow$ Europe stays tight & $-$ \\
US load-queue wait & $n$-weighted mean & US cannot export $\Rightarrow$ HH softens, Europe starved & $-$ \\
net export pressure & $n$-weighted mean & $z$(loadings) $-$ $z$(load queue); both legs point the same way & $+$ \\
\bottomrule
\end{tabular}
\end{center}

Volumes are in m$^3$ of LNG; loading and discharge rates are amortised m$^3$ per day
across berth hours so that a closed visit integrates to one cargo; ages and waits are
in days and hours respectively. The project's signal catalogue documents the
remaining twenty-six keys, the two bases, and the confidence columns (dispersion,
open fraction, estimated fraction) that accompany every value.

\paragraph{A recorded inconsistency.} The catalogue's own text gave net export
pressure two incompatible signs in two places. The positive sign was adopted on the
economics above and recorded as a prior; the data returned a negative point estimate
at both horizons, at $|t|\le1.6$, which neither confirms nor falsifies it.
